\chapter{Introduction}

This chapter introduces the motivation, problem, research questions, scope, and claimed contributions of the thesis.

\section{Background and Context}

Knowledge graphs have become fundamental infrastructure for organising and reasoning about complex information across domains ranging from web search to biomedical research \citep{kgquality2024}. They are also increasingly relevant to retrieval-augmented generation, including GraphRAG approaches that use graph structure to organise retrieved evidence and support more context-aware generation \citep{edge2024graphrag,han2025graphrag}. Traditional knowledge graphs often represent relationships between entities, such as a person working at a company or a company being located in a city. Listing \ref{lst:entity-centred-example} records both facts, but it does not represent the company's founding as an event. It therefore does not state when that event happened, where it happened, or who took part in it.

\begin{lstlisting}[language={},numbers=none,caption={A small entity-centred RDF example.},label={lst:entity-centred-example}]
@prefix ex: <http://example.org/> .

ex:Founder ex:worksAt ex:Company .
ex:Company ex:locatedIn ex:CapeTown .
\end{lstlisting}

Event-centric knowledge graphs (EKGs) address this limitation by placing events at the centre of representation \citep{sem2011}. Listing \ref{lst:event-centred-example} extends the earlier example by representing the founding as a resource with participants, a place, and a date. This supports questions about what happened, when and where it happened, and who was involved. This thesis focuses on evaluating the quality of that event representation rather than every possible reasoning task that could be performed over it.

\noindent\begin{minipage}{0.96\textwidth}
\begin{lstlisting}[language={},numbers=none,caption={An event-centred extension of the earlier RDF example.},label={lst:event-centred-example}]
@prefix ex: <http://example.org/> .
@prefix sem: <http://semanticweb.cs.vu.nl/2009/11/sem/> .
@prefix xsd: <http://www.w3.org/2001/XMLSchema#> .

ex:CompanyFounding a sem:Event ;
    sem:hasActor ex:Founder, ex:Company ;
    sem:hasPlace ex:CapeTown ;
    sem:hasBeginTimeStamp "2012-05-10"^^xsd:date .
\end{lstlisting}
\end{minipage}

The shift from entity-centric to event-centric knowledge representation introduces significant complexity. Events can have temporal extent, participants, locations, external identifiers, source provenance, and relationships with other events. Several EKG systems have emerged to address different aspects of this challenge. EventKG integrates event information from multiple sources, including Wikipedia, Wikidata, and YAGO, among others, and provides temporal and spatial context for large numbers of events \citep{eventkg2018}. ASER focuses on extracting eventuality knowledge from text, capturing relationships between activities and states \citep{aser2019}, and the Open Event Knowledge Graph (OEKG) emphasises multilingual event representation and cross-lingual linking \citep{oekg2023}.

As EKG construction methods have developed, a gap has emerged in how their quality is evaluated. Current evaluation approaches tend to focus on narrow aspects, such as extraction precision and recall, temporal coverage, or downstream task performance, without providing a broad view of EKG quality \citep{kgquality2021}. This fragmentation makes it difficult to compare EKG systems, identify their strengths and weaknesses across multiple dimensions, or guide targeted improvements. For example, EventKG reports event and attribute coverage, including time and location information, whereas ASER reports extraction accuracy and performance on commonsense reasoning tasks \citep{eventkg2018,aser2019}. These results answer different questions and cannot be read as the same measure of quality.

The need for broader intrinsic evaluation extends beyond academic comparison. Organisations deploying EKGs for decision support need evidence about temporal information, structural linkage, duplicate candidates, event descriptions, and external links. These indicators do not prove factual correctness or task fitness, but they expose limitations that a narrow extraction or coverage score can hide.

\section{Problem Statement}

Current EKG evaluation approaches usually assess one part of quality at a time and report results in different forms. Extraction studies may report precision, recall, and F1 against annotated reference data; resource studies may report event or attribute coverage; and downstream studies may report task accuracy \citep{eventkg2018,aser2019,chronographer2023}. These studies count different things and divide their results by different populations. Their values therefore answer different questions and cannot be directly averaged.

Bringing these evaluations together does not mean reducing them to one score. It means defining a common profile that keeps each measure separate while making its population, denominator, and interpretation clear. Without such a profile, temporal, structural, descriptive, and linking weaknesses can remain hidden behind a narrow result.

The core research problem is therefore the lack of a common, reproducible way to assess several intrinsic properties of an RDF-based EKG together. Within the scope of the literature review in Chapter 2, no identified framework brings these concerns into one RDF-EKG evaluation process. This thesis investigates which dimensions belong in that profile, how they can be measured, and what is learned when they are considered together.

\section{Research Questions, Aim, and Objectives}

This section defines the questions used to evaluate whether the work meets its purpose, then states the aim and breaks that aim into concrete objectives.

\subsection{Research Questions}

The research addresses three primary questions:

\begin{enumerate}
    \item \textbf{RQ1:} What dimensions of quality are relevant for event-centric knowledge graphs, and how can they be systematically measured?
    \item \textbf{RQ2:} What quality differences and trade-offs are revealed when the dimensions identified in RQ1 are evaluated together rather than through selected narrow approaches?
    \item \textbf{RQ3:} What quality patterns emerge when the framework is applied to controlled datasets and independently produced public RDF event graphs?
\end{enumerate}

RQ1 identifies and formalises the dimensions needed for the common profile. RQ2 tests whether bringing those dimensions together addresses the fragmented view described in the problem statement. RQ3 examines whether the resulting profile supports useful interpretation across controlled and independently produced graphs.

\subsection{Research Aim and Objectives}

The aim of this research is to investigate EKG quality requirements and metrics in order to design, formalise, implement, and evaluate a broader intrinsic framework for RDF-based EKGs.

To achieve this aim, the research pursues five objectives:

\begin{enumerate}
    \item Design a conceptual evaluation framework that identifies and organises the relevant dimensions of EKG quality.
    \item Formalise a set of computable intrinsic measures for EKGs by combining established KG quality measures with EKG-specific measures such as temporal density, entity richness, external mapping coverage, and predicate usage patterns.
    \item Implement a reproducible, MIT-licensed evaluation toolkit that operationalises the framework through SPARQL-based analysis and graph-theoretic measures.
    \item Evaluate the framework through a controlled sensitivity experiment on high-quality, mixed-quality, and low-quality synthetic RDF datasets, and test execution feasibility and descriptive applicability on public RDF event-graph artefacts.
    \item Compare the multidimensional framework with narrow single-dimension views and selected general RDF/KG quality tools to show what quality differences and trade-offs would otherwise be missed.
\end{enumerate}

The comparison in Objective 5 tests whether the full profile exposes defects and trade-offs hidden by selected metric subsets and general RDF/KG quality tools. The profile can inform hypotheses about timeline analysis, integration, exploratory querying, or graph analysis, but this thesis does not test whether its intrinsic values predict downstream task performance.

\section{Scope}

This research is limited to event-centric knowledge graphs represented in RDF. It uses OWL vocabulary terms and SPARQL queries where required \citep{rdf11concepts2014,owl2overview2012,sparql11query2013}. The default profile assumes direct SEM event typing and SEM-style temporal properties. It evaluates nine selected intrinsic dimensions: graph connectivity and structure, redundancy and duplication, temporal consistency, minimal event-profile alignment, completeness, closed-profile type alignment, entity richness, external mapping coverage, and predicate usage.

The empirical evaluation uses three controlled synthetic RDF datasets designed to represent high-, mixed-, and low-quality conditions and two independently produced public RDF cases: the OEKG event layer and three ChronoGrapher graph artefacts. Across the synthetic generation profiles, the probabilities of labels, dates, places, descriptions, relations, and external links are reduced, while label reuse, coarse dates, and type mismatches are increased. These datasets test sensitivity to known changes; because the changes target measured properties, they are not independent construct validation.

The complete OEKG event layer tests whether the implementation can process a large published EKG on the available hardware. Its 13 published event-layer files were copied into a cleaned working folder, where seven invalid literal escapes were repaired without adding semantic facts. The smaller ChronoGrapher artefacts test the same metrics on a different public RDF representation. Chapter 3 explains the generation, cleaning, and preparation procedures, while Appendix \ref{app:oekg_run_details} records the main commands and hashes.

The framework evaluates EKG data rather than proposing new event extraction, relation detection, or graph construction methods. It also measures information quality rather than query-engine performance or the behaviour of downstream machine-learning systems. Domain-specific criteria, multilingual evaluation, downstream task utility, and causal correctness are therefore treated as extensions. Causal links are important for some EKG applications, but evaluating whether one event genuinely caused another requires domain assumptions, annotation, or downstream reasoning tasks; it is not claimed as a completed metric dimension here.

\section{Significance and Contributions}

This research addresses a practical and methodological gap in EKG evaluation. For the research community, it provides a systematic way to compare EKG quality across multiple dimensions rather than relying on isolated metrics. For practitioners, it provides a command-line tool that computes the declared measures without manual calculation and identifies concrete weaknesses such as missing dates, sparse event descriptions, low external-link coverage, high redundancy, or fragmented graph structure.

This research makes four primary contributions:

\begin{enumerate}
    \item A conceptual evaluation framework identifying nine selected intrinsic dimensions of RDF-EKG quality.
    \item A set of 28 distinct metric calculations across the nine dimensions. The software reports them through 32 named fields because four fields repeat a result needed by another reporting module.
    \item A command-line implementation that computes the metrics using SPARQL queries and graph analysis.
    \item Empirical evidence from controlled synthetic profiles, the OEKG event layer, and public ChronoGrapher artefacts, demonstrating controlled sensitivity, real-graph execution feasibility, and additional diagnostic coverage beyond narrow metric subsets and four general RDF/KG quality tools.
\end{enumerate}

The first three contributions are developed in Chapter 3. That chapter traces the nine dimensions to the literature, defines the 28 calculations and their 32 named result fields, and describes the command-line implementation. Chapter 4 supplies the empirical evidence through the controlled synthetic experiment, execution on OEKG and ChronoGrapher, and comparison with narrow views and four general RDF/KG tools. Chapter 5 discusses what this evidence supports and where the framework remains limited.

\section{Thesis Structure}

Chapter 2 reviews EKGs, temporal evaluation, KG quality, RDF validation, and related evaluation frameworks, and identifies the gaps that motivate the proposed framework. Chapter 3 explains how the literature informed the framework design, how the metrics were formalised, and how the implementation was verified. Chapter 4 reports the controlled synthetic results, the general-tool comparison, and the OEKG and ChronoGrapher runs. Chapter 5 interprets those results, answers the research questions, and states the limitations. Chapter 6 concludes the thesis and identifies priority future work.
